\documentclass[11pt]{article}
\usepackage[margin=1in]{geometry}
\usepackage{amsmath,amssymb,booktabs}
\usepackage[hidelinks]{hyperref}
\setlength{\parskip}{5pt}
\setlength{\parindent}{0pt}

\newcommand{\eone}{\eta_{D_1}}
\newcommand{\etwo}{\eta_{D_2}}
\newcommand{\ecomp}{\eta_C}

\title{A Stationary Theory of Degree-4 Circuit Competition in a Two-Block ReLU Network}
\author{}
\date{}

\begin{document}
\maketitle

\begin{abstract}
The frozen-representation theory compares a direct \(1+1+1+1\to4\) gate in block 2 with a
compositional \(2+2\to4\) gate in block 2.  When block 1 is also allowed to learn the degree-4
target, there is an additional direct route.  A restricted comparison of one compositional
route with two coherently pooled direct routes moves the equal-norm boundary from
\(Q_2/Q_1=1/\sqrt8\) to \(Q_2/Q_1=1/2\).  This is not, however, the unrestricted global
minimum: the block-1 direct route can share the readout with either kind of block-2 gate.  In
that full stationary problem it is a common branch and cancels from the choice between the two
block-2 mechanisms.  The block-2 gate boundary remains \(Q_2/Q_1^{(1)}=1/\sqrt8\), but the
winning solution is generically a hybrid containing both block-1 direct signal and block-2
signal.  Thus a pure \(2+2\to4\) phase cannot be expected without suppressing, penalizing, or
geometrically misaligning the block-1 direct route.  We give the general competition equation,
unequal-decay and finite-temperature extensions, a stationary order-parameter formulation for
joint \(d2+d4\) training, and controlled experiments that distinguish the alternatives.
\end{abstract}

\section{Architecture and three candidate gates}

Write the two-block residual network as
\begin{align}
h_1(x)&=h_0(x)+\sigma(V_1h_0(x)),\\
h_2(x)&=h_1(x)+\sigma(V_2h_1(x)),\\
\widehat y_4(x)&=w^\top h_2(x),
\end{align}
where \(\sigma\) is ReLU.  Let
\begin{itemize}
    \item \(Q_1^{(0)}\) be the norm of each relevant singleton Walsh vector in \(h_0\);
    \item \(Q_1^{(1)}\) be the norm of each relevant singleton Walsh vector in \(h_1\);
    \item \(Q_2\) be the norm of each of the two parent vectors \(q_{12},q_{34}\) in \(h_1\).
\end{itemize}
The three ideal balanced gates are
\begin{align}
D_1 &: 1+1+1+1\longrightarrow4 \quad\hbox{in block 1},\\
D_2 &: 1+1+1+1\longrightarrow4 \quad\hbox{in block 2},\\
C &: 2+2\longrightarrow4 \quad\hbox{in block 2}.
\end{align}
Using \(|\kappa_4|=1/4\) and \(|\kappa_2|=1/2\), their target-mode efficiencies per hidden
weight norm are
\begin{equation}
\boxed{
\eone=\frac{Q_1^{(0)}}8,
\qquad
\etwo=\frac{Q_1^{(1)}}8,
\qquad
\ecomp=\frac{Q_2}{2\sqrt2}.}
\label{eq:etas}
\end{equation}
This section initially treats \(Q_1^{(0)},Q_1^{(1)},Q_2\) as fixed order parameters.  Their
self-consistent adaptation is considered later.

\section{One homogeneous route: why weight decay compares efficiencies}

For one route of efficiency \(\eta\), let \(u\) denote its hidden-weight norm and \(v\) its
readout norm.  An aligned rank-one realization of target amplitude \(A\geq0\) obeys
\begin{equation}
A=\eta uv.
\end{equation}
The smallest quadratic parameter cost is
\begin{equation}
\min_{\eta uv=A}\frac12(u^2+v^2)=\frac{A}{\eta},
\label{eq:one-route-cost}
\end{equation}
attained at \(u=v=\sqrt{A/\eta}\).  With squared error and common weight decay, the reduced
objective is
\begin{equation}
\mathcal J_\eta(A)=\frac12(1-A)^2+\lambda\frac{A}{\eta}.
\end{equation}
Consequently
\begin{equation}
A_\eta^*=\left[1-\frac{\lambda}{\eta}\right]_+,
\qquad
\mathcal J_\eta^*=
\begin{cases}
\dfrac{\lambda}{\eta}-\dfrac{\lambda^2}{2\eta^2},&\lambda<\eta,\\[4pt]
\dfrac12,&\lambda\geq\eta.
\end{cases}
\label{eq:amplitude}
\end{equation}
Thus larger \(\eta\) has lower stationary cost for every nonsuppressed solution.  A common
\(\lambda\) sharpens selection and sets the suppression scale, but does not move a boundary
defined by equality of efficiencies.  Different normalizations of the MSE rescale the stated
suppression value, not the circuit boundary.

\section{Parallel routes and the role of feature alignment}

The new issue is that the block-1 and block-2 target features are added in the same residual
stream before a shared readout.  Let their normalized target-feature directions have cosine
\(\rho\).  If their individual efficiencies are \(a,b\), their best joint efficiency is the
largest singular value of the two-column feature map:
\begin{equation}
\eta_{\parallel}^2(a,b;\rho)=\frac12\left[a^2+b^2+
\sqrt{(a^2-b^2)^2+4a^2b^2\rho^2}\right].
\label{eq:parallel-general}
\end{equation}
Two limits are useful:
\begin{equation}
\eta_{\parallel}(a,b;0)=\max(a,b),
\qquad
\eta_{\parallel}(a,b;1)=\sqrt{a^2+b^2}.
\label{eq:parallel-limits}
\end{equation}
Orthogonal routes cannot share readout norm and the more efficient one wins.  Coherently aligned
routes share the readout and add in root-sum-square.  Since the network can in principle place a
block-1 target feature in the same residual coordinate as a block-2 target feature, the aligned
limit is the natural zero-temperature global-minimum prediction unless auxiliary constraints
prevent alignment.  Measuring \(\rho\) is therefore essential.

\section{Restricted pooled-direct boundary}

First impose a restricted two-family comparison:
\begin{itemize}
    \item the direct family may use \(D_1\) and \(D_2\) together;
    \item the compositional family may use \(C\), but is forbidden from also using \(D_1\).
\end{itemize}
Then
\begin{equation}
\eta_D=\eta_{\parallel}(\eone,\etwo;\rho_D),
\qquad
\eta_C=\ecomp,
\end{equation}
and the restricted phase boundary is
\begin{equation}
\boxed{\ecomp=\eta_{\parallel}(\eone,\etwo;\rho_D).}
\label{eq:restricted-general}
\end{equation}
In the fully aligned case this becomes
\begin{equation}
\frac{Q_2}{2\sqrt2}
=\frac18\sqrt{\left(Q_1^{(0)}\right)^2+\left(Q_1^{(1)}\right)^2},
\end{equation}
or
\begin{equation}
\boxed{
Q_{2,c}=\frac{1}{2\sqrt2}
\sqrt{\left(Q_1^{(0)}\right)^2+\left(Q_1^{(1)}\right)^2}.}
\label{eq:restricted-boundary}
\end{equation}
If the singleton norms are equal, \(Q_1^{(0)}=Q_1^{(1)}=Q_1\), this gives the simple new
number
\begin{equation}
\boxed{\left.\frac{Q_2}{Q_1}\right|_{c,\mathrm{restricted}}=\frac12.}
\label{eq:half}
\end{equation}
The shift \(1/\sqrt8\to1/2\) is a real prediction for this restricted comparison: two coherent
direct production sites lower the direct family's norm cost.  It is also an upper endpoint of
the alignment-dependent boundary.  For equal direct efficiencies, varying
\(0\leq|\rho_D|\leq1\) moves the threshold continuously from \(1/\sqrt8\) to \(1/2\).

\section{The unrestricted two-block minimum is hybrid}

Equation~\eqref{eq:half} is not the generic phase boundary of the unrestricted architecture.
The reason is simple: if \(D_1\) is available, it can cooperate with a compositional block-2
gate just as it can cooperate with a direct block-2 gate.  The two correct competing families
are therefore
\begin{align}
H_D&=D_1+D_2, &
\eta_{H_D}&=\eta_{\parallel}(\eone,\etwo;\rho_D),\\
H_C&=D_1+C, &
\eta_{H_C}&=\eta_{\parallel}(\eone,\ecomp;\rho_C).
\end{align}
The general stationary competition equation is
\begin{equation}
\boxed{
\eta_{\parallel}(\eone,\ecomp;\rho_C)
=
\eta_{\parallel}(\eone,\etwo;\rho_D).}
\label{eq:full-boundary}
\end{equation}
If both families can attain the same alignment---in particular, if both attain coherent
alignment---\(\eta_{\parallel}\) is monotone in its second argument, so the common \(D_1\)
branch cancels.  The choice of block-2 gate is then
\begin{equation}
\ecomp\gtrless\etwo,
\end{equation}
and hence
\begin{equation}
\boxed{
\left.\frac{Q_2}{Q_1^{(1)}}\right|_{c,\mathrm{block\ 2\ gate}}
=\frac1{\sqrt8}.}
\label{eq:old-boundary}
\end{equation}
The additional direct route does not move this particular boundary.  Instead, it changes the
nature and cost of both phases.  In the aligned limit,
\begin{equation}
\eta_{H_D}=\sqrt{\eone^2+\etwo^2},
\qquad
\eta_{H_C}=\sqrt{\eone^2+\ecomp^2},
\end{equation}
so both tolerate stronger decay than their single-block counterparts.

Moreover, the optimal \(H_C\) solution is not a pure compositional circuit.  At the balanced
minimum, the two target-amplitude contributions are proportional to squared efficiencies:
\begin{equation}
f_{D_1\mid H_C}=\frac{\eone^2}{\eone^2+\ecomp^2},
\qquad
f_{C\mid H_C}=\frac{\ecomp^2}{\eone^2+\ecomp^2}
=\frac{8Q_2^2}{(Q_1^{(0)})^2+8Q_2^2}.
\label{eq:hybrid-fraction}
\end{equation}
For \(Q_1^{(0)}=Q_1^{(1)}\), the compositional share is \(1/2\) at the block-2 gate boundary
and grows smoothly above it.  More generally, achieving \(f_C\geq p\) requires
\begin{equation}
\frac{Q_2}{Q_1^{(0)}}\geq
\frac1{2\sqrt2}\sqrt{\frac{p}{1-p}}.
\label{eq:purity}
\end{equation}
This is a purity condition, not a phase boundary.  It quantifies how strong the parents must be
to make the compositional contribution dominate the unavoidable block-1 shortcut.

\section{Unequal decay exposes how to force composition}

Let the weight-decay coefficients be \(\lambda\rho_1\) on target-specific block-1 weights,
\(\lambda\rho_2\) on block-2 weights, and \(\lambda\rho_w\) on the readout.  In the aligned
case, a route in block \(\ell\) contributes squared effective efficiency
\(\eta_\ell^2/(\rho_\ell\rho_w)\).  The restricted direct-versus-compositional boundary is
\begin{equation}
\frac{\ecomp^2}{\rho_2}
=\frac{\eone^2}{\rho_1}+\frac{\etwo^2}{\rho_2},
\end{equation}
or
\begin{equation}
\boxed{
8Q_2^2=(Q_1^{(1)})^2+\frac{\rho_2}{\rho_1}(Q_1^{(0)})^2.}
\label{eq:unequal-decay}
\end{equation}
Penalizing target formation in block 1 more strongly (\(\rho_1\gg\rho_2\)) continuously
returns the old \(1/\sqrt8\) boundary.  Protecting block 1 from decay
(\(\rho_1\to0\)) does the opposite: weight decay alone can never disfavor the untaxed direct
shortcut.  This distinguishes ``protect the already learned parent representation'' from
``leave every block-1 direction unregularized.''  A useful setup protects the parent-producing
subspace while still penalizing new block-1 \(d4\) directions.

For the unrestricted hybrid competition, the common aligned \(D_1\) term again cancels:
\begin{equation}
\frac{\eone^2}{\rho_1}+\frac{\ecomp^2}{\rho_2}
\gtrless
\frac{\eone^2}{\rho_1}+\frac{\etwo^2}{\rho_2}.
\end{equation}
Thus unequal block-1 decay changes the shortcut fraction and suppression scale but, under this
fixed-representation closure, not which gate block 2 prefers.

\section{A stationary theory for simultaneous parent and target learning}

When \(V_1\) is trainable, \(Q_1^{(0)},Q_1^{(1)},Q_2\) are not externally fixed.  Substituting
their initial or final measured values into a frozen-representation boundary is not a complete
theory.  A dynamics-free generalization is nevertheless possible.  Let
\(\Phi(Q_1^{(0)},Q_1^{(1)},Q_2)\) be the minimum loss plus regularization needed to realize the
singleton representation and the two supervised degree-2 parents.  For
\(r\in\{H_D,H_C\}\), define the stationary free energy
\begin{equation}
\boxed{
F_r=\min_{Q_1^{(0)},Q_1^{(1)},Q_2,\,A\geq0}
\left\{
\Phi(Q_1^{(0)},Q_1^{(1)},Q_2)
+\frac12(1-A)^2
+\lambda\frac{A}{\eta_r(Q_1^{(0)},Q_1^{(1)},Q_2)}
\right\}.}
\label{eq:landau}
\end{equation}
The joint-training phase boundary is \(F_{H_C}=F_{H_D}\).  This is a small order-parameter
theory rather than a parameter-space dynamical equation.  It automatically allows the
compositional phase to strengthen \(Q_2\), and global decay to weaken it, if those changes lower
the total stationary cost.

A toy calculation shows the feedback explicitly.  Let
\(\eta_P=Q_1^{(0)}/(2\sqrt2)\) be the efficiency for producing each degree-2 parent in block 1,
let both parent targets have amplitude \(B\), and let their common hidden norm be \(Q_2\).
After minimizing their readout directions, the regularization cost of the two parents is
\begin{equation}
\Phi_{\rm toy}(Q_2)=\frac{Q_2^2}{\eta_P^2}+\frac{B^2}{Q_2^2}.
\label{eq:parent-potential}
\end{equation}
Without downstream use, its minimizer is
\begin{equation}
Q_{2,D}^4=B^2\eta_P^2.
\end{equation}
For a pure compositional downstream route of fixed amplitude \(A\), the additional cost is
\(2\sqrt2 A/Q_2\), and the self-consistent parent strength obeys
\begin{equation}
\boxed{
\frac{Q_{2,C}^4}{\eta_P^2}-B^2-\sqrt2 A Q_{2,C}=0.}
\label{eq:self-consistent-q2}
\end{equation}
Hence \(Q_{2,C}>Q_{2,D}\): downstream composition creates a stationary incentive to amplify
its own parents.  This positive feedback is absent from the frozen theory.  Equation
\eqref{eq:landau}, with a more faithful empirically or analytically derived \(\Phi\), is the
minimal route to a genuine two-block phase diagram without writing training dynamics.

The fixed-representation results remain a controlled limit of
Eq.~\eqref{eq:landau}: take \(\Phi\) infinitely stiff around prescribed values of the three
order parameters.  A sensible progression is therefore stiff parents \(\to\) softly pinned
parents \(\to\) fully joint \(d2+d4\) training.

\section{Connection to a Boltzmann or field-theory description}

For explicitly isotropic Langevin dynamics satisfying detailed balance, the stationary
parameter distribution is
\begin{equation}
p(\theta)\propto
\exp\left[-\frac{1}{T}\left(
\mathcal L(\theta)+\frac\lambda2\|\theta\|^2\right)\right].
\end{equation}
The zero-temperature saddle is exactly the minimum-cost theory above.  Integrating parameters
at fixed order parameters \(Q=(Q_1^{(0)},Q_1^{(1)},Q_2,A,\ldots)\) defines an effective action
\begin{equation}
\Gamma(Q)=E_{\min}(Q)-T S(Q),
\end{equation}
whose leading saddle is Eq.~\eqref{eq:landau}.  At one loop, a circuit family \(r\) acquires
the usual curvature and multiplicity terms
\begin{equation}
F_r(T)\simeq E_r^*+\frac{T}{2}\log\det{}'H_r-T\log g_r.
\label{eq:one-loop}
\end{equation}
In the crude exclusive-route approximation, two equivalent direct placements give
\(g_D=2\), \(g_C=1\), and equal probability occurs at
\begin{equation}
E_C^*=E_D^*-T\log2.
\label{eq:entropy-shift}
\end{equation}
Thus finite temperature favors the more numerous direct family even when the zero-temperature
costs are equal.  Neural networks have much larger continuous degeneracies, so \(T\log2\) is
only an illustration; the Hessian/entropy term, not route counting alone, is the controlled
object.  Ordinary SGD with momentum need not sample this Boltzmann law.  The connection should
therefore be tested with deliberately injected Langevin noise, after the zero-temperature
stationary predictions are established.

\section{Controlled experiments implied by the theory}

No simultaneous-training experiment is needed first.  The cleanest next test keeps the parent
order parameters fixed while adding only the previously missing target-specific block-1
degrees of freedom.

\paragraph{Representation.}
Construct or pretrain a base stream with four singleton vectors and the two parent vectors,
then freeze the parent-producing subspace.  Add a fresh trainable block-1 target adapter whose
input contains the four singletons.  This adapter can make \(D_1\) but cannot alter \(Q_2\).
Train it together with block 2 and the \(d4\) readout.  Applying decay to the adapter, block 2,
and readout preserves the minimum-norm comparison while protecting the learned parents.

\paragraph{Axes.}
Sweep measured \(Q_2/Q_1^{(1)}\) densely across both \(1/\sqrt8\simeq0.354\) and \(1/2\), vary
\(Q_1^{(0)}/Q_1^{(1)}\), and use a decay range large enough to approach the stationary
minimum without suppressing the target.  A second sweep of \(\rho_1/\rho_2\) tests
Eq.~\eqref{eq:unequal-decay}.

\paragraph{Necessary measurements.}
Separately measure the degree-4 Walsh vectors
\begin{equation}
q_4^{(1)}=\widehat{h_1}(\{1,2,3,4\}),
\qquad
\Delta q_4^{(2)}=\widehat{\sigma(V_2h_1)}(\{1,2,3,4\}),
\end{equation}
their readout contributions, and their alignment
\begin{equation}
\rho=\frac{q_4^{(1)}\cdot\Delta q_4^{(2)}}
{\|q_4^{(1)}\|\,\|\Delta q_4^{(2)}\|}.
\end{equation}
Classifying only block-2 preactivations cannot distinguish a pure compositional circuit from
the predicted \(D_1+C\) hybrid.  Record target-specific parameter norms as well, so that the
cost \(A/\eta_{\rm eff}\) can be checked directly.

\paragraph{Four decisive ablations.}
\begin{center}
\begin{tabular}{@{}lll@{}}
\toprule
Allowed routes & Main prediction & Purpose\\
\midrule
\(D_2,C\) & gate boundary \(1/\sqrt8\) & reproduce frozen baseline\\
\(D_1,D_2\) & \(\eta_D=\eta_{\parallel}(\eone,\etwo;\rho_D)\) & test pooling/alignment\\
\(D_1,C\) & hybrid fraction Eq.~\eqref{eq:hybrid-fraction} & test shared-readout synergy\\
\(D_1,D_2,C\) & block-2 boundary Eq.~\eqref{eq:full-boundary} & test full theory\\
\bottomrule
\end{tabular}
\end{center}
The restricted \(1/2\) boundary should be tested by an ablation that really excludes \(D_1\)
from the compositional family; otherwise the unrestricted hybrid theory predicts the block-2
transition at \(1/\sqrt8\), not \(1/2\).

Only after these fixed-order-parameter tests should the parent stiffness be relaxed.  Add a
quadratic penalty or auxiliary loss that softly pins \(Q_2\), sweep its stiffness, and fit the
effective potential \(\Phi\) in Eq.~\eqref{eq:landau}.  Fully simultaneous \(d2+d4\) training
is the zero-stiffness endpoint.  This sequence separates three effects that are otherwise
confounded: extra circuit topology, representation adaptation, and training-time lock-in.

\section{Main conclusions}

\begin{enumerate}
    \item A restricted pure-composition versus pooled-direct calculation has the new aligned,
    equal-singleton boundary \(Q_2/Q_1=1/2\), with the general form
    Eq.~\eqref{eq:restricted-boundary}.

    \item In the unrestricted architecture, block-1 direct signal can cooperate with either
    block-2 gate.  The correct competition is \(D_1+C\) versus \(D_1+D_2\),
    Eq.~\eqref{eq:full-boundary}.  With comparable alignment, the block-2 gate boundary remains
    \(Q_2/Q_1^{(1)}=1/\sqrt8\).

    \item The compositional side is generically a hybrid, not a pure phase.  For aligned
    routes, its compositional target share is given by Eq.~\eqref{eq:hybrid-fraction}.  Forcing
    a nearly pure \(2+2\to4\) circuit requires suppressing the block-1 target shortcut, applying
    stronger decay to new block-1 target directions, or making \(Q_2/Q_1^{(0)}\) sufficiently
    large.

    \item Joint training can still be treated statically by minimizing an effective free energy
    over \(Q_1^{(0)},Q_1^{(1)},Q_2\) and target amplitude.  This stationary order-parameter
    theory is the natural zero-temperature limit of a later Boltzmann/field-theory treatment.
\end{enumerate}

\end{document}
